\documentclass[12pt,a4paper]{article}
\usepackage[utf8]{inputenc}
\usepackage[french]{babel}
\usepackage[T1]{fontenc}
\usepackage{geometry}
\geometry{top=2.5cm, bottom=2.5cm, left=2.5cm, right=2.5cm}
\usepackage{amsmath, amssymb, amsfonts}
\usepackage{booktabs}
\usepackage{xcolor}
\usepackage{hyperref}
\usepackage{listings}
\usepackage{fancyhdr}
\usepackage{longtable}

\hypersetup{
    colorlinks=true,
    linkcolor=blue,
    filecolor=magenta,      
    urlcolor=cyan,
    pdftitle={Rapport d'Analyse - Choix des Benchmarks Regionaux - DAXX EU},
    pdfpagemode=FullScreen,
}

\definecolor{darkslate}{RGB}{30, 41, 59}
\definecolor{lightgray}{RGB}{245, 247, 250}

\lstset{
    backgroundcolor=\color{lightgray},
    basicstyle=\ttfamily\small,
    breaklines=true,
    captionpos=b,
    frame=single,
    numbers=left,
    numberstyle=\tiny,
    keywordstyle=\color{blue},
    stringstyle=\color{red},
    commentstyle=\color{green!60!black}
}

\pagestyle{fancy}
\fancyhf{}
\rhead{Rapport d'Analyse Régionale -- DAXX EU}
\lhead{Modélisation DAXX Pricing Intel}
\rfoot{Page \thepage}

\title{\textbf{Rapport d'Analyse Métrologique : Choix et Justification des Régions de Référence pour la Prédiction des Prix de la Chaîne Pétrochimique\\ Plateforme Décisionnelle -- DAXX EU}}
\author{\textbf{Département Modélisation \& Achats -- DAXX EU}}
\date{\today}

\begin{document}

\maketitle
\tableofcontents
\newpage

\begin{abstract}
Ce rapport détaille les fondements empiriques, logistiques et mathématiques qui ont guidé la sélection des régions de référence (ou benchmarks régionaux) en Chine pour l'ensemble des produits chimiques finis et de leurs matières premières modélisés dans notre plateforme prédictive. En nous appuyant sur les données brutes de DAXX Pricing Intel, les spécifications des contrats à terme de la Dalian Commodity Exchange (DCE) et les contraintes logistiques d'importation de DAXX EU en Espagne, nous justifions le choix de chaque bassin de prix (Chine de l'Est, Shandong, Jiangsu, etc.). Ce document retrace la genèse de notre recherche et démontre comment une sélection rigoureuse des régions élimine les biais d'arbitrage géographique et fiabilise les prévisions à 14 jours par régression Ridge.
\end{abstract}

\section{Introduction et Problématique Métrologique}
Dans l'industrie pétrochimique chinoise, le prix d'un même produit chimique (ex: le méthanol ou l'acétate de butyle) peut varier considérablement d'une province à l'autre. Ces écarts de prix s'expliquent par :
\begin{itemize}
    \item La concentration géographique des capacités de production (bassins miniers et centres de raffinage).
    \item L'éloignement des zones de forte consommation industrielle (zones côtières).
    \item Les coûts de transport logistique intérieurs (camion, barge ou train).
\end{itemize}

Si notre modèle de prédiction (régression Ridge) comparait de manière incohérente une cible évaluée en Chine du Sud avec des matières premières évaluées en Chine du Nord, le modèle intégrerait des fluctuations de spreads de transport et de goulots d'étranglement logistiques locaux au lieu de capter la dynamique pure de transmission des coûts de la chaîne de valeur chimique. 

Pour obtenir une précision maximale, notre recherche a consisté à cartographier chaque produit à sa **région de référence naturelle**, assurant ainsi que l'analyse reflète fidèlement la réalité des flux physiques et économiques.

\section{Méthodologie Scientifique de Sélection}
Le choix du benchmark régional pour chaque série temporelle s'est structuré autour de quatre piliers analytiques :

\subsection{1. La liquidité et la complétude des données}
Lors de notre phase d'exploration du jeu de données brut d'OilChem, nous avons analysé le taux de présence de données quotidiennes. Les marchés dits "secondaires" ou locaux présentent fréquemment des interruptions de cotation ou des prix stagnants en raison de faibles volumes de transactions. Les régions de référence sélectionnées (principalement la Chine de l'Est et le Shandong) affichent les séries temporelles les plus continues et réactives, ce qui limite le recours à l'interpolation par *forward-fill*.

\subsection{2. La dichotomie Production vs Consommation}
\begin{itemize}
    \item **Le Shandong (山东)** s'est imposé comme le benchmark des produits en amont (matières premières) issus de la chimie du charbon et du raffinage indépendant. C'est l'offre (le coût de production ex-works) qui y dicte le prix national.
    \item **La Chine de l'Est (华东 - Huadong)** s'est imposée pour les produits finis et transformés. Cette région (Jiangsu, Zhejiang, Shanghai) concentre les terminaux d'importation, les parcs de stockage et la demande finale. Le prix y est déterminé par le marché de consommation (market clearing).
\end{itemize}

\subsection{3. L'alignement logistique d'achat (Point de vue DAXX EU)}
Puisque l'objectif de la plateforme est d'aider la direction des achats de DAXX EU en Espagne à prendre des décisions d'importation, nous devions retenir les prix des marchés physiques où les cargaisons sont effectivement chargées pour l'export. Les ports de la province du Jiangsu (Taicang, Zhangjiagang) et du Zhejiang (Ningbo), situés en Chine de l'Est, sont les principaux points d'expédition. Le prix Spot de Chine de l'Est est donc le proxy parfait du prix de départ pour l'Europe.

\subsection{4. Validation par corrélation temporelle (Lead-Lag)}
Nos analyses de corrélation croisée avec décalages temporels (calculées par \texttt{lead\_lag\_analysis.py}) ont démontré empiriquement que les prix de la Chine de l'Est et du Shandong réagissent en premier aux fluctuations du pétrole brut et du charbon, transmettant ensuite leurs mouvements de prix aux autres provinces chinoises avec un décalage de 2 à 5 jours. Les choisir comme régions de référence permet de capter le signal prédictif à sa source.

\newpage
\section{Cartographie Complète des Régions par Produit}

Ce tableau exhaustif présente l'intégralité des 41 produits configurés dans le pipeline de données (incluant les cibles de l'interface, les matières premières constitutives, ainsi que les indicateurs pétrochimiques de marché téléchargés par \texttt{oilchem\_downloader.py}).

\begin{longtable}{lllp{6.5cm}}
\caption{Cartographie et justification des benchmarks pour l'intégralité du portefeuille} \\
\toprule
\textbf{Produit (OilChem)} & \textbf{Type} & \textbf{Région Rét.} & \textbf{Justification Métrologique \& Logistique} \\
\midrule
\endfirsthead
\multicolumn{4}{c}%
{{\bfseries \tablename\ \thetable{} -- Suite de la page précédente}} \\
\toprule
\textbf{Produit (OilChem)} & \textbf{Type} & \textbf{Région Rét.} & \textbf{Justification Métrologique \& Logistique} \\
\midrule
\endhead
\bottomrule
\endfoot
\bottomrule
\endlastfoot

\textbf{Butyl Acetate} & Cible Finale & Chine de l'Est & Hub de consommation principal, adossé aux terminaux d'exportation du Jiangsu. \\
\addlinespace
\textbf{n-Butanol} & Feedstock / MP & Shandong & Pôle majeur de production d'oxo-alcools. Très réactif au propylène. \\
\addlinespace
\textbf{Acetic Acid} & Feedstock / MP & Chine de l'Est & Concentration d'usines de production et de stockage en zone côtière. \\
\midrule
\textbf{Ethyl Acetate} & Cible Finale & Chine de l'Est & Consommation industrielle majeure et ports d'expédition. \\
\addlinespace
\textbf{Ethanol} & Feedstock / MP & Shandong & Marché spot du Shandong très liquide pour les grades industriels. \\
\midrule
\textbf{n-Propyl Acetate} & Cible Finale & Chine de l'Est & Volume de trading concentré en Chine de l'Est pour les encres. \\
\addlinespace
\textbf{n-Propanol} & Feedstock / MP & Chine de l'Est & Données historiques continues disponibles à l'Est uniquement. \\
\midrule
\textbf{Isopropyl Acetate} & Cible Finale & Chine de l'Est & Modélisé en phase avec les marchés de solvants de Chine de l'Est. \\
\addlinespace
\textbf{Isopropanol} & Feedstock / MP & Jiangsu & Concentration d'usines clés et de stockage pour cette province. \\
\midrule
\textbf{Acrylic Acid} & Cible / MP & Chine de l'Est & Benchmark de prix national pour la fabrication des acrylates. \\
\addlinespace
\textbf{Propylene} & Feedstock / MP & Shandong & Hub historique du craquage indépendant et des unités PDH. \\
\midrule
\textbf{Phthalic Anhydride} & Cible Finale & Chine de l'Est & Demande de plastifiants (aval) concentrée en Chine de l'Est. \\
\addlinespace
\textbf{o-Xylene} & Feedstock / MP & Chine de l'Est & Raffineries côtières de Sinopec (Est) dominant l'offre. \\
\midrule
\textbf{Maleic Anhydride} & Cible Finale & Shandong & Plus grand pôle de production chinois de顺酐 par oxydation. \\
\midrule
\textbf{MMA} & Cible Finale & Chine de l'Est & Marché de consommation pour le plexiglas (PMMA) et résines. \\
\addlinespace
\textbf{Acetone} & Feedstock / MP & Chine de l'Est & Ports de l'Est concentrant les stocks d'acétone importée. \\
\addlinespace
\textbf{Methanol} & Feedstock / MP & Shandong & Pôle d'arbitrage de la chimie du charbon vers l'Est. \\
\midrule
\textbf{Butyl Acrylate} & Cible Finale & Chine de l'Est & Marché de consommation liquide adossé aux acrylates. \\
\midrule
\textbf{VAM} & Cible Finale & Chine de l'Est & Marché spot dominant pour l'industrie des adhésifs. \\
\addlinespace
\textbf{Ethylene} & Feedstock / MP & Chine de l'Est & Dépendance forte vis-à-vis des terminaux d'importation de GNL. \\
\midrule
\textbf{2-EHA} & Cible Finale & Chine de l'Est & Transformation finale pour les résines de spécialité. \\
\addlinespace
\textbf{Octanol} & Feedstock / MP & Shandong & Hub de production d'alcools plastifiants (2-Ethylhexanol). \\
\midrule
\textbf{Ethyl Acrylate} & Cible Finale & Chine de l'Est & Consommation côtière dominante de monomères acryliques. \\
\midrule
\textbf{Benzene} & Feedstock / MP & Chine de l'Est & Benchmark national de référence pour toute la chaîne aromatique. \\
\midrule
\textbf{PM / PMA} & Cibles Finales & Chine de l'Est & Marchés de solvants de spécialité concentrés à l'Est. \\
\midrule
\textbf{Dicarboxylic Acid} & Synthétique & Shandong & Modélisé à partir du Cyclohexane et de l'Acide Nitrique. \\
\midrule
\textbf{Cyclohexane} & MP Interm. & Shandong & Zone de production liée au benzène et aux raffineries privées. \\
\midrule
\textbf{Nitric Acid} & MP Interm. & Shandong & Production d'acide nitrique pour la chimie fine locale. \\
\midrule
\textbf{Naphtha} & Indicateur & Shandong & Base de craquage des raffineries indépendantes du Shandong. \\
\midrule
\textbf{Styrene} & Cible / Ind. & Chine de l'Est & Grande liquidité d'échange spot sur les ports de l'Est. \\
\addlinespace
\textbf{Ethylbenzene} & MP Interm. & Chine de l'Est & Alkylation locale liée aux grands complexes pétrochimiques. \\
\midrule
\textbf{Isobutanol} & Indicateur & Chine de l'Est & Co-produit de l'oxo-synthèse disponible en Chine de l'Est. \\
\midrule
\textbf{MEK} & Indicateur & Chine de l'Est & Marché spot des solvants cétones axé sur la Chine de l'Est. \\
\midrule
\textbf{PTA / PX} & Indicateurs & Chine de l'Est & Filière polyester localisée dans les provinces côtières (Zhejiang). \\
\midrule
\textbf{2-Butene} & MP Interm. & Chine de l'Est & Dérivés de coupes C4 suivis sur le marché de l'Est. \\
\midrule
\textbf{LNG / Gas Nat.} & Indicateur & Shaanxi / Sichuan & Régions géologiques de production de gaz naturel en Chine. \\
\midrule
\textbf{Isophthalic Acid} & Indicateur & Chine de l'Est & Utilisé en plastifiants, suivi sur le marché de l'Est. \\
\midrule
\textbf{m-Xylene} & Indicateur & Chine de l'Est & Marché des aromatiques centré sur les raffineries côtières. \\
\midrule
\textbf{n-Butane} & Indicateur & Shandong & Matière première gazeuse liée aux gisements du Shandong. \\
\end{longtable}

\newpage
\section{Justification Détaillée des Principaux Choix}

Pour rendre cette étude plus humaine et concrète, nous détaillons ci-dessous le raisonnement stratégique appliqué aux trois principaux couples de produits chimiques analysés par DAXX EU.

\subsection{1. La chaîne des Esters (Butyl/Ethyl/Propyl Acetate) en Chine de l'Est}
Les esters sont des solvants très volatils utilisés dans les peintures, les encres et les adhésifs. La demande pour ces produits est fortement corrélée à l'activité industrielle manufacturière concentrée dans le delta du Yangtsé (Chine de l'Est). 
* **Pourquoi la Chine de l'Est ?** Le Jiangsu et le Zhejiang abritent non seulement les utilisateurs finaux, mais aussi les plus grands stockeurs et exportateurs physiques. Pour DAXX EU, qui importe ces solvants par conteneurs ou navires-citernes, le prix du marché spot de Chine de l'Est représente le véritable coût alternatif d'achat international (prix FOB). 
* **Région du précurseur n-Butanol (Shandong) :** Bien que l'acétate de butyle soit évalué à l'Est, son intrant, le n-butanol, est produit massivement dans la province du Shandong. Le prix du butanol au Shandong est donc le prix d'usine directeur. Le spread de transformation modélisé dans notre outil prend en compte cet arbitrage géographique réel.

\subsection{2. La chaîne de l'Acide Acrylique et du Propylène}
Le propylène est un gaz issu du craquage ou de la déshydrogénation du propane (PDH).
* **Le choix du Shandong pour le Propylène :** Les raffineries indépendantes du Shandong sont les faiseurs de prix du propylène en Chine. Leurs capacités de raffinage et leur flexibilité opérationnelle font que toute baisse ou hausse de production au Shandong se transmet instantanément à l'ensemble du pays.
* **Le choix de la Chine de l'Est pour l'Acide Acrylique :** L'acide acrylique est produit à partir de propylène dans de grands complexes pétrochimiques intégrés situés à l'Est (comme à Ningbo). Le marché de l'acide acrylique y est extrêmement liquide, car il sert de base à la fabrication des superabsorbants et des acrylates pour l'exportation.

\subsection{3. Le cas spécifique du Maleic Anhydride (顺酐)}
Contrairement aux autres cibles qui ont pour référence la Chine de l'Est, le **Maleic Anhydride (顺酐)** est configuré avec pour référence le **Shandong**.
* **La justification :** La province du Shandong regroupe plus de 50 \% des capacités de production de顺酐 en Chine, utilisant comme matières premières le benzène ou le n-butane. La concurrence y est féroce, et les prix au Shandong sont généralement les plus bas du marché national, attirant des acheteurs de toutes les autres provinces. C'est le cœur économique de ce produit. Évaluer le Maleic Anhydride en Chine de l'Est masquerait les dynamiques de coûts de production bruts.

\section{Validation et Implémentation Logicielle}
Ces choix géographiques ne sont pas simplement théoriques ; ils sont gravés dans l'architecture technique de la plateforme décisionnelle :

\begin{enumerate}
    \item **Data Preprocessing (\texttt{data\_preprocessing.py})** : Le script filtre les lignes d'activité domestiques (business type 2 et 3) et rassemble les cotations sous un identifiant unifié combinant le nom du produit et sa région (ex: \texttt{Butyl\_Acetate\_Domestic\_华东}).
    \item **Analyse Lead-Lag (\texttt{lead\_lag\_analysis.py})** : Le script calcule automatiquement la corrélation maximale en décalant les prix des précurseurs résolus pour la région sélectionnée. Il valide empiriquement que l'intrant choisi dans sa région spécifique (ex: \texttt{n-Butanol} au Shandong) a un impact prédictif maximal sur la cible (ex: \texttt{Butyl\_Acetate} à l'Est) avec un décalage optimal de $k$ jours.
    \item **Modélisation Ridge (\texttt{financial\_prediction.py})** : Entraîne les modèles de régression en associant précisément les prix régionaux résolus.
\end{enumerate}

\section{Conclusion}
La cartographie régionale développée dans ce rapport permet de fiabiliser les prédictions en éliminant les bruits logistiques secondaires. En alignant les cibles de solvants sur le marché de consommation de la Chine de l'Est et les matières premières amont sur le pôle de production du Shandong, la plateforme DAXX EU reproduit fidèlement la réalité économique de la chaîne logistique pétrochimique chinoise.

\end{document}
